%% ===================================================================
%%  DoomTex -- one pdflatex run renders one frame of Doom.
%%
%%      input.txt  +  state.tex  ->  [ pdflatex doom.tex ]
%%                                        ->  doom.pdf  +  state.tex
%%
%%  Run it again to advance one more frame.  Everything -- the fixed
%%  point maths, the raycaster, the textures, the AI -- is TeX macro
%%  expansion.  No Lua, no shell escape, no external tools.
%%
%%  See README.md.  For a keyboard-driven loop, use ./play.sh
%% ===================================================================
\documentclass{article}
\usepackage[paperwidth=340bp,paperheight=326bp,margin=10bp]{geometry}
\usepackage{doomtex}
\pagestyle{empty}
\setlength{\parindent}{0pt}

\begin{document}
\ExplSyntaxOn

%% ---- boot -------------------------------------------------------
%% Both tables are computed by TeX on the first run and cached to disk,
%% so every later frame pays nothing for them.
\dtexLoadSinTable:
\dtexLoadTextures:

%% ---- read the player's intent, and the world as we left it ------
\dtexReadInput
\dtexLoadState

%% ---- simulate exactly one tick ----------------------------------
\dtexTick

%% ---- render ------------------------------------------------------
\dtexCastAll
\vbox
  {
    \offinterlineskip
    \dtexRenderAudio
    \kern 6bp
    %% The viewport: the column-stack frame, then the sprite overlay,
    %% then the weapon, all layered at the same origin.
    \hbox
      {
        \rlap { \dtexRenderFrame }
        \rlap { \dtexDrawSprites }
        \rlap { \dtexDrawWeapon }
        \rlap { \dtexDamageTint }
        \kern \dim_eval:n { \c_dtex_screen_w_int \c_dtex_px_dim }
      }
    \kern 8bp
    \dtexRenderHUD
  }

%% ---- hand the world to the next run ------------------------------
\dtexSaveState
\ExplSyntaxOff
\end{document}
